\section{Appendix: Liturgical Resolution}

The First Sunday of Advent, Year B (PC-S01-B), falls on 29 November 2026 in the United States national calendar. The liturgical year is 2027 and the color is violet. The Sunday Mass has one form, with no separate vigil set. The Gloria is omitted and the Creed is said. The independent weekday cycle is I. The \textit{Roman Missal}, Third Edition for the United States (2011), supplies the common First Sunday prayers and antiphons for the Sunday cycles. The U.S. \textit{Lectionary for Mass}, second typical edition, vol.~I, no.~2, supplies the Year B readings. Its Latin predecessor in the order of readings, \textit{Ordo lectionum Missae} (1981), explains the choice of a Gospel about the Lord's final coming for the First Sunday of Advent (\textit{Praenotanda} 93).

\begin{branchtable}
\bid{entrance-chant-choice} & U.S. General Instruction of the Roman Missal (GIRM) 48 & Permitted & Entrance & Missal antiphon or an authorized entrance chant. The local choice is unknown.\\
\bid{sprinkling} & GIRM 51, Sunday Mass & Permitted & Introductory rites & The blessing and sprinkling of water can replace the Penitential Act.\\
\bid{psalm-sung-choice} & GIRM 61, Lectionary 173--174 & Permitted & After Isaiah & Psalm 80 on the Sunday route. Seasonal responses and psalms are identified below. The Gradual route is also permitted.\\
\bid{acclamation-chant-choice} & GIRM 62--63 & Permitted & Before Gospel & The Alleluia verse in no.~2 or the permitted Graduale choice. Local music is unknown.\\
\bid{advent-wreath-within-mass} & U.S. Book of Blessings 1509, 1517--1520 & Permitted & Close of the intercessions & The presider uses 1519 or 1520, then the first candle is lit. No local celebration is known.\\
\bid{preparation-chant-choice} & GIRM 74 & Permitted & Preparation of gifts & No particular chant is appointed in the Sunday's Missal proper.\\
\bid{preface-eucharistic-prayer-choice} & Sunday rubric, Order of Mass 33, GIRM 365 & Coupled choice & Eucharistic Prayer & The Sunday directs to Preface I of Advent. The choice of Eucharistic Prayer must respect its Preface provisions. Prayer IV has an invariable Preface.\\
\bid{communion-chant-choice} & GIRM 87 & Permitted & Communion & Missal antiphon or an authorized Communion chant. The local choice is unknown.\\
\bid{concluding-blessing-choice} & Sunday rubric, Solemn Blessings I.1 & Permitted & Dismissal & The Advent solemn blessing may replace the ordinary blessing.\\
\end{branchtable}

The Sunday has no Sequence, alternative reading length or separate proper Offertory antiphon. The Advent wreath has a different purpose and a different place from the solemn blessing at the end of Mass. The U.S. ritual permits the wreath blessing on this Sunday or the evening before it. The order for the household is a separate provision (no.~1514).

\subsection{When an Advent common psalm is sung}

The Lectionary provides common responses and psalms which, when sung, may replace the corresponding texts so that the congregation can participate more readily (GIRM 61). No.~173 gives an Advent response asking the Lord to come and set his people free. No verse number accompanies it. No.~174 gives two psalms with their own responses. The following are the underlying scriptural passages in Challoner, with the response passage first. The authorized musical texts have their own phrasing. The U.S. headings give whole verses. The Latin order specifies how the strophes begin and end.

\paragraph{Psalm 25 (24).} The response draws on v.~1. The three strophes are 4--5, 8--9, and 10 with 14. The Latin order specifies 4bc--5ab, 8--9, 10 and 14. In Challoner the first strophe begins after the clause concerning the unjust in 24:4 and ends at the identification of God as Savior in 24:5.

\begin{quote}
To thee, O Lord, have I lifted up my soul.

Shew, O Lord, thy ways to me, and teach me thy paths. Direct me in thy truth, and teach me; for thou art God my Saviour;

The Lord is sweet and righteous: therefore he will give a law to sinners in the way. He will guide the mild in judgment: he will teach the meek his ways.

All the ways of the Lord are mercy and truth, to them that seek after his covenant and his testimonies. The Lord is a firmament to them that fear him: and his covenant shall be made manifest to them.
\end{quote}

Augustine, on this psalm, pictures people who need pardon before they can set foot on God's road. They must also be willing to be taught once they are on it. He understands the mercy and truth of God's ways through the two advents, pardon extended and deeds judged. Bellarmine notices the difference between being shown a road and being given help to walk it. To the penitent who fears being lost again he points out the promise of instruction to those who fear the Lord. When the psalm replaces the Sunday strophes, a plea to be taught follows Isaiah's confession of wandering.\footnote{Augustine, \textit{Exposition on Psalm} 24, §§4--5, 8--10, 14. Bellarmine, \textit{Commentary on the Book of Psalms}, Psalm 24, explanations 4--5, 8--10, 14.}

\paragraph{Psalm 85 (84).} The response draws on the opening request in v.~8. The three strophes are 9--10, 11--12 and 13--14. The Latin order begins 9ab--10. The U.S. first strophe hears God's proclamation of peace and passes at once to the nearness of salvation. The additional recipients listed in the full Challoner v.~9 are outside the sung selection.

\begin{quote}
Shew us, O Lord, thy mercy;

I will hear what the Lord God will speak in me: for he will speak peace unto his people: Surely his salvation is near to them that fear him: that glory may dwell in our land.

Mercy and truth have met each other: justice and peace have kissed. Truth is sprung out of the earth: and justice hath looked down from heaven.

For the Lord will give goodness: and our earth shall yield her fruit. Justice shall walk before him: and shall set his steps in the way.
\end{quote}

This option includes both the verse about truth growing from the earth and the verse about its fruit. Augustine hears the prophet turning from the surrounding noise to hear God speak peace. Bellarmine asks who will bring such peace. He sees the Savior in the prayer for salvation and attends at the close to the justice by whose light people can find their way. The longer psalm at this point of Mass lets the worshippers follow the prayer from listening to taking a step. Neither common psalm changes the acclamation before the Gospel or the Communion antiphon printed in the Missal.\footnote{Augustine, \textit{Exposition on Psalm} 84, §§6--13. Bellarmine, \textit{Commentary on the Book of Psalms}, Psalm 84, explanations 7--13.}

\section{Appendix: Scope and Qualifications}

This study concerns the national observance. It follows the Missal antiphons and the appointed response from Psalm 80. No information establishes which music, Eucharistic Prayer, wreath rite or final blessing an individual church chose. The printed biblical English belongs to the public-domain Challoner revision of the Douay--Rheims Bible in the Gutenberg 1581 edition. The approved U.S. Lectionary controls proclamation. The adaptations of biblical verses in the Missal and the protected wording of the Lectionary, its response, the English ritual prayers and modern Latin Missal prayers are omitted. The Scripture and the comments here are for study.

The U.S. Missal from a third-party digital copy gives the Sunday on p.~139, the Preface on p.~534 and the Advent blessing on p.~674. A U.S. publisher sample independently supplies the blessing with different pagination. The bishops' conference supplies the three orations with commentary and the current conclusion to the Collect. The official 2008 list of corrections to the Latin third typical edition records no change to the Advent I texts considered here. The Latin digital copy is of the earlier 2002 text. A complete 2008 altar book was unavailable for direct collation. The common Advent I Missal record, the Year B appointment record and the national calendar determination distinguish those sources.

Jerome's and Thomas's Latin commentaries give the wording they knew. The Vulgate-based study text and the U.S. proclamation differ at several points in Isaiah. Bede's Latin Gospel joins prayer to watching at Mark 13:33, as Challoner does. An English heading for Augustine's psalm exposition may give a different psalm number from the Latin title, and editions number verses differently according to how they count the superscription. O'Sullivan's historical English version of Bellarmine omits some philological material. Neither the Latin original behind that abridgment nor an independent Greek text of Chrysostom or Cyril serves as the basis for a textual reconstruction here. The Fathers and later doctors expound their respective biblical passages. The bishops' conference gives direct commentary on the modern orations. The connections drawn between those prayers and the separate biblical expositions belong to the present study.

\section{References}
\begin{enumerate}
\item \textit{Douay--Rheims Bible}, Challoner revision, Project Gutenberg 1581. Biblical study English, with adaptations and native verse boundaries identified at the appointed passages.
\item USCCB, \href{https://bible.usccb.org/bible/readings/112926.cfm}{First Sunday of Advent, 29 November 2026}, Lectionary 2, and \href{https://www.usccb.org/resources/2027cal.pdf}{U.S. Liturgical Calendar for 2027}, printed pp.~5, 11. National observance, readings, response and acclamation.
\item \textit{Roman Missal}, Third Edition, United States (2011), First Sunday of Advent, Order of Mass 33 and Solemn Blessings I.1. \href{https://archive.org/download/Roman-Missal-Third-Edition/Roman\%20Missal.pdf}{Digital U.S. altar book}, \href{https://b2c-cbp-assets.s3.amazonaws.com/pdfs/9780899420738.pdf}{Catholic Book Publishing sample} and \href{https://www.liturgyoffice.org.uk/Missal/Music/Antiphonary.pdf}{ICEL Antiphonary}, First Sunday. The \href{https://media.musicasacra.com/books/latin_missal2002.pdf}{Latin third-typical-edition digital reproduction} and variation list in \textit{Notitiae} 44 (2008), pp.~367--387, distinguish the 2002 text and the emended reprint.
\item \textit{Ordo lectionum Missae}, second typical edition (1981), \textit{Praenotanda} 93 and nos.~1--2, 173--174. U.S. General Instruction of the Roman Missal, \href{https://www.usccb.org/prayer-and-worship/the-mass/general-instruction-of-the-roman-missal/girm-chapter-2}{chapter II} and chapter VII, no.~365, on the rites and their permitted choices.
\item USCCB, \href{https://www.usccb.org/prayer-and-worship/liturgical-year-and-calendar/advent/commentary-on-advent-proper-prayers}{\textit{Commentary on the Proper Prayers of Advent from the Roman Missal}}, First Sunday, and Committee on Divine Worship Newsletter, February 2021, pp.~7--8, on the Collect conclusion.
\item U.S. \textit{Book of Blessings} (Liturgical Press, 1989), chapter on the blessing of an Advent wreath, nos.~1509, 1514, 1517--1520. Archdiocese of Louisville, \href{https://www.archlou.org/wp-content/uploads/2023/09/Advent-2023.pdf}{\textit{Season of Advent 2023}}, pp.~3--4, on the wreath and sung psalms. \href{https://www.modernpsalter.com/us/Lectionary.aspx?n=174\&s=0}{Modern Psalter, Common Advent Psalms}, U.S. Lectionary no.~174, settings 1--2, for the U.S. passage boundaries.
\item Jerome, \href{https://www.documentacatholicaomnia.eu/03d/0347-0420,_Hieronymus,_In_Isaiam_Pars_4,_LT.pdf}{\textit{Commentary on Isaiah}}, book XVII, Latin part IV, pp.~149--159. The people's accusation, the assurance in the Latin text and the clay's appeal to the Father.
\item Augustine of Hippo, \textit{Expositions on the Psalms}, Latin Psalms 24, 79 and 84, historical NPNF I.8 and the respective New Advent texts headed \href{https://www.newadvent.org/fathers/1801025.htm}{Psalm 25}, \href{https://www.newadvent.org/fathers/1801080.htm}{Psalm 80} and \href{https://www.newadvent.org/fathers/1801085.htm}{Psalm 85}. The Church's trust, the vine, confession, and the earth's fruit. \href{https://www.augustinus.it/latino/lettere/lettera_204_testo.htm}{\textit{Letter} 199 to Hesychius}, I.2--3, Latin, on the people included in the command to watch.
\item Thomas Aquinas, \href{https://www.corpusthomisticum.org/cis50.html}{\textit{Expositio super Isaiam ad litteram}}, chapter 64, §85024, and \href{https://www.corpusthomisticum.org/c1c.html}{\textit{Super I ad Corinthios}}, chapter 1, lecture 1, §86234, on defective acts of justice and gifts shared in charity.
\item John Chrysostom, \textit{Homilies on First Corinthians}, \href{https://www.newadvent.org/fathers/220101.htm}{Homily 1} and \href{https://www.newadvent.org/fathers/220102.htm}{Homily 2}. The recipients of grace, praise for the gifts and the warning implicit in the promise.
\item Bede, \textit{In Marci evangelium expositio}, on Mark 13:33--37, Latin Wikisource transcription from Corpus Corporum, citing PL 92, on the departure of the master, the porter's task and the universal command.
\item Robert Bellarmine, \textit{A Commentary on the Book of Psalms}, trans. John O'Sullivan (Dublin, 1866), Psalms 24, 79 and 84. Confidence in God's promises, the husbandman of the vine and the Virgin's fruit.
\item Cyril of Jerusalem, \href{https://www.newadvent.org/fathers/310115.htm}{\textit{Catechetical Lecture} XV}, §§1--3. The two comings in the Creed and the identity of the judge.
\end{enumerate}
